\section{Mathematical analysis}\label{sec:mechanisms}

The following elementary results explain the search geometry and the conditions
needed to trust a proposal. They do not establish convergence of the hybrid
solver. Coordinates in the derivations are real-valued, macro dimensions and
orientations are fixed, and touching boundaries are legal. The implementation
uses floating-point coordinates and separate endpoint validation.

\subsection{Feasible regions and coordinate slices}\label{sec:geometry-analysis}

Let $\mathcal B$ contain all coordinate bounds and fixed-coordinate equalities.
Let $\mathcal E$ be the unordered hard-macro pairs, with
$m=|\mathcal E|$, and let $H_{e,b}$ be the closed halfspace corresponding to
branch $b\in\{1,2,3,4\}$ of Equation~\eqref{eq:disjunction} for pair $e$.

\begin{proposition}[Finite polyhedral decomposition]\label{prop:geometry}
The feasible set is a union of at most $4^m$ compact convex polyhedra:
\begin{equation}
 \mathcal F
 =\mathcal B\cap\bigcap_{e\in\mathcal E}\bigcup_{b=1}^4 H_{e,b}
 =\bigcup_{\sigma\in\{1,2,3,4\}^{\mathcal E}} P_\sigma,
 \qquad
 P_\sigma=\mathcal B\cap\bigcap_{e\in\mathcal E}H_{e,\sigma(e)}.
 \label{eq:polyhedral-union}
\end{equation}
The polyhedra may overlap or be empty.
\end{proposition}
\begin{proof}
Distributing the finite intersections over unions selects one non-overlap
inequality for each pair. Each resulting set is an intersection of closed
halfspaces and fixed-coordinate equalities, hence a closed convex polyhedron.
The canvas bounds make it bounded, and therefore compact.
\end{proof}

A coordinate slice need not itself be an interval. Fix a legal placement and
vary only $x_i$ for a movable hard macro $i$. Write $\mathcal H$ for the hard
macros, and define the perpendicular blockers and required horizontal separation
by
\begin{equation}
 S_i=\left\{j\in\mathcal H\setminus\{i\}:
 |y_i-y_j|<\frac{h_i+h_j}{2}\right\},
 \qquad d_{ij}=\frac{w_i+w_j}{2}.
\end{equation}
The exact legal slice is
\begin{equation}
 I_i=\left[\frac{w_i}{2},W-\frac{w_i}{2}\right]
 \setminus\bigcup_{j\in S_i}(x_j-d_{ij},x_j+d_{ij}).
 \label{eq:coordinate-slice}
\end{equation}
Indeed, a nonblocker already satisfies a vertical separation inequality; a
blocker excludes precisely the open horizontal interval in
Equation~\eqref{eq:coordinate-slice}. Its endpoints remain legal because contact
is allowed. Thus $I_i$ is a finite union of closed intervals, including possible
singletons. The $y_i$ slice follows by exchanging axes; a soft macro has only
its canvas interval, and a fixed macro has a singleton slice.

The routine \code{legal\_axis\_range} does not enumerate this union. It bounds
one interval using the current left/right relations to perpendicular blockers
and adds a numerical clearance. This is a conservative search restriction;
subsequent perpendicular moves can change those relations. It does not imply
that all CD sweeps remain in one $P_\sigma$. The smooth stage can move several
coordinates through overlapping intermediate states, but its legalized output
is passed to CD without comparison against the initial SDF score. Geometry
alone therefore supplies no monotonicity guarantee for the complete pipeline.

\subsection{Incremental consistency and acceptance error}\label{sec:incremental-analysis}

Let $\Psi(p)$ denote the cache obtained by recomputing all components at placement
$p$. For net $n$, let $a_n(p)$ be its weighted HPWL contribution, $d_i(p)$ the
density-bin vector of macro $i$, $r_n^o(p)$ the raw routing vector of net $n$, and
$b_i^o(p)$ the hard-macro routing contribution, for orientation $o\in\{H,V\}$.
Soft macros have $b_i^o=0$. The required aggregate invariant is
\begin{equation}
 T=\sum_n a_n(p),\qquad D=\sum_i d_i(p),\qquad
 R^o=\sum_n r_n^o(p),\qquad M^o=\sum_i b_i^o(p),
 \label{eq:cache-invariant}
\end{equation}
with cached pin positions and net extrema also equal to their values at $p$.
Normalization, smoothing, and tail aggregation then produce the three scores.

For a move $U_i$ of macro $i$, only its pins, its incident nets, and its own
geometric contributions change. If these dependencies are complete and each
contribution is evaluated correctly, subtracting the old terms and adding the
new ones gives
\begin{equation}
 \widehat U_i\Psi(p)=\Psi(U_i p),
 \label{eq:cache-commutation}
\end{equation}
where $\widehat U_i$ is the corresponding cache update. This follows term by
term from Equation~\eqref{eq:cache-invariant}; induction preserves the invariant
over a sequence of moves. Reversing the substitutions restores the prior state
in exact arithmetic. The result is conditional on complete dependency tracking,
correct extrema updates, and correct grid contributions. Floating-point
subtract/add operations need not cancel, and their summation order can differ
from full recomputation.

\begin{proposition}[Acceptance under bounded score error]\label{prop:acceptance}
Let $S_p$ and $S_q$ be the incremental states at candidate endpoints $p$ and $q$.
Suppose
\begin{equation}
 |\widehat J(S_p)-\proxy(p)|\le\epsilon,\qquad
 |\widehat J(S_q)-\proxy(q)|\le\epsilon.
\end{equation}
If $\widehat J(S_q)\le\widehat J(S_p)-\tau$, then
\begin{equation}
 \proxy(q)-\proxy(p)\le-\tau+2\epsilon.
 \label{eq:acceptance-error}
\end{equation}
In particular, $\tau>2\epsilon$ guarantees a strict canonical decrease.
\end{proposition}
\begin{proof}
Add and subtract the incremental endpoint values:
\begin{align*}
 \proxy(q)-\proxy(p)
 &=\bigl[\proxy(q)-\widehat J(S_q)\bigr]
   +\bigl[\widehat J(S_q)-\widehat J(S_p)\bigr]
   +\bigl[\widehat J(S_p)-\proxy(p)\bigr]\\
 &\le\epsilon-\tau+\epsilon. \qedhere
\end{align*}
\end{proof}

Absolute component bounds $\epsilon_{\rm wl}$, $\epsilon_{\rm den}$, and
$\epsilon_{\rm cong}$ imply
$\epsilon=\epsilon_{\rm wl}+\tfrac12\epsilon_{\rm den}
+\tfrac12\epsilon_{\rm cong}$ by Equation~\eqref{eq:canonical} and the triangle
inequality. The measured relative discrepancies in Section~\ref{sec:evaluator}
are samples on prescribed traces, not uniform absolute bounds. They therefore
do not certify every CD acceptance at its $10^{-9}$ threshold or every
assignment-batch acceptance at $10^{-7}$. Full endpoint rescoring checks that
endpoint; it does not reconstruct or certify all preceding decisions.

\subsection{Tail aggregation and the cost of a move}\label{sec:tail-analysis}

For $z\in\mathbb R^G$ and a fixed integer $1\le k\le G$, let
$z_{(1)}\ge\cdots\ge z_{(G)}$ denote its sorted entries. Define the top-$k$ mean
$A_k(z)=k^{-1}\sum_{j=1}^k z_{(j)}$.

\begin{lemma}[Tail representation and stability]\label{lem:tail}
With $(u)_+=\max\{u,0\}$,
\begin{equation}
 A_k(z)
 =\max_{\substack{0\le w_j\le1/k\\\sum_j w_j=1}}w^\top z
 =\min_{\theta\in\mathbb R}\left\{\theta+
 \frac1k\sum_{j=1}^G(z_j-\theta)_+\right\}.
 \label{eq:tail-representation}
\end{equation}
It is convex and piecewise affine, and for any $\delta\in\mathbb R^G$,
\begin{equation}
 |A_k(z+\delta)-A_k(z)|
 \le\min\left\{\|\delta\|_\infty,
          \frac{\|\delta\|_2}{\sqrt{k}},
          \frac{\|\delta\|_1}{k}\right\}.
 \label{eq:tail-stability}
\end{equation}
\end{lemma}
\begin{proof}
The maximizing weights put mass $1/k$ on $k$ largest entries. For any $\theta$,
$\sum_{j=1}^k z_{(j)}\le k\theta+\sum_j(z_j-\theta)_+$, with equality at
$\theta=z_{(k)}$, proving the threshold form. The maximum is equivalently over
finitely many $k$-element subsets, so it is convex and piecewise affine. For any
admissible $w$, its norms satisfy $\|w\|_1=1$,
$\|w\|_2\le1/\sqrt{k}$, and $\|w\|_\infty\le1/k$.
Bounding the difference between the two maxima by
$\max_w|w^\top\delta|$ and applying H\"older's inequality gives all three bounds.
\end{proof}

For a grid with $G\ge10$ cells and nonnegative normalized density loads $\rho$,
the canonical density rule is
$C_{\rm den}=\tfrac12 A_{\lfloor0.1G\rfloor}(\rho)$. The implementation drops
zero cells but retains the denominator $\lfloor0.1G\rfloor$; this is equivalent
to padding with zeros when fewer cells are occupied. For $G<10$, its separate
occupied-cell mean rule requires separate treatment. Congestion applies
$A_{\lfloor0.05(2G)\rfloor}$ to the concatenated horizontal and vertical
normalized, smoothed, and macro-adjusted loads; when that index is zero, the
code uses their maximum. Selection ties can make these functions
nondifferentiable. Convexity in grid loads does not imply convexity in placement
coordinates, because the load maps themselves depend on the geometry.

Local cache mutation also does not imply local total scoring cost. Let
$T_{\rm mut}(i)$ be the work to update macro $i$'s pins, incident nets, and
changed grid contributions, and let $T_{\rm sel}(G,k)$ be the cost of the
implemented top-$k$ selection, with $T_{\rm sel}(0,0)=0$. For $G\ge10$, set
$k=\lfloor0.1G\rfloor$ and let $n_{\rm occ}\le G$ count the nonzero density
cells. The public \code{move} call returns the current score and has the structure
\begin{equation}
 T_{\rm move}(i)=T_{\rm mut}(i)+O(G)
 +T_{\rm sel}(n_{\rm occ},\min\{k,n_{\rm occ}\})
 +T_{\rm sel}(2G,k).
 \label{eq:move-cost}
\end{equation}
Density normalization, congestion smoothing, and aggregation traverse the grid.
The present implementation therefore retains at least linear work in $G$ per
scored move; it does not maintain an incremental order-statistic structure.
The microbenchmarks measure the savings from localized component updates at the
reported grid sizes, not a runtime bound depending only on macro degree.

\Needspace{18\baselineskip}
\subsection{What negative curvature guarantees}\label{sec:curvature}

A negative-curvature direction can give descent even when the gradient is
nonzero. The guarantee concerns a specified smooth function and an unmodified
step; feasibility and subsequent legalization impose additional conditions.
All norms in the following proposition are Euclidean for vectors and the
induced operator norm for matrices.

\begin{proposition}[Signed curvature decrease]\label{prop:curvature}
Let $f$ be twice continuously differentiable near $p$, and let $v$ be a unit
vector with $v^\top\nabla^2 f(p)v\le-\eta<0$. Choose
$s\in\{-1,1\}$ such that $s\nabla f(p)^\top v\le0$. Fix $t>0$ and suppose $f$ is twice
continuously differentiable on an open neighborhood of the segment
$\{p+usv:0\le u\le t\}$, and its Hessian is $L$-Lipschitz there. Then
\begin{equation}
 f(p+tsv)-f(p)\le-\frac{\eta}{2}t^2+\frac{L}{6}t^3.
 \label{eq:curvature-taylor}
\end{equation}
For $L>0$ and $0<t\le3\eta/(2L)$, this is at most $-\eta t^2/4$.
For $L=0$, the first bound is $-\eta t^2/2$ throughout the assumed segment.
\end{proposition}
\begin{proof}
Taylor's theorem with Lipschitz Hessian gives a remainder bounded in absolute
value by $L\|tsv\|_2^3/6=Lt^3/6$. The linear term
$ts\nabla f(p)^\top v$ is nonpositive, and the quadratic term is at most
$-\eta t^2/2$. The stated step bound makes $Lt^3/6\le\eta t^2/4$.
\end{proof}

The gradient need not vanish. However, the selected sign must admit a feasible
step for this to yield feasible descent. For example, $f(x)=x-x^2$ on
$[0,1/4]$ has a strict constrained minimum at $x=0$ despite $f'(0)=1$ and
$f''(0)=-2$. The favorable sign points outside the interval. Projecting or
legalizing a step changes its endpoint, so Equation~\eqref{eq:curvature-taylor}
no longer bounds the resulting value without an additional displacement bound.

Eigensolver accuracy is a separate issue. For a fixed real symmetric operator
$\mathsf H$, a unit vector $v$, Rayleigh quotient $\rho=v^\top\mathsf H v$,
and residual $r=\mathsf H v-\rho v$,
\begin{equation}
 \operatorname{dist}(\rho,\operatorname{spec}\mathsf H)\le\|r\|_2.
 \label{eq:eigen-residual}
\end{equation}
To see this, expand $v=\sum_j a_j u_j$ in an orthonormal eigenbasis. Then
$\|r\|_2^2=\sum_j a_j^2(\lambda_j-\rho)^2
\ge\min_j|\lambda_j-\rho|^2$, since $\sum_j a_j^2=1$.
A small residual locates some eigenvalue, not necessarily the smallest one;
$\rho<0$ itself implies $\lambda_{\min}(\mathsf H)<0$ by the Rayleigh bound.
Neither statement validates the operator as the Hessian of the scored function.

Lane B applies automatic differentiation to a legacy proxy $J_B$, with rectangle
intersections, clamps, and top-$k$ selection. A classical Hessian need not exist
at their branch boundaries. In Section~\ref{sec:curvature-diagnostic}, guarded
float64 probes remove canvas-branch changes but retain disagreement in the
piecewise density and congestion terms. Those checks do not validate a classical
Hessian of the composite proxy at the measured endpoint, despite the small
same-operator eigensolver residual.
Moreover, Lane B legalizes and polishes its proposals, then accepts by the
canonical score rather than $J_B$. Proposition~\ref{prop:curvature} consequently
explains a proposal rationale; its assumptions have not been established for
that execution. A matched comparison with random and shuffled directions must
measure canonical improvement retained after legalization and polish.

\subsection{Assignment interactions}\label{sec:assignment-analysis}

Let $p^0$ be the frozen background and $p^{i\leftarrow j}$ the placement obtained
by moving only macro $i$ to slot $j$. Using an exact scalar objective $F$ to
separate modeling error from cache error, define
$\Delta_{ij}=F(p^{i\leftarrow j})-F(p^0)$. The assignment proposal solves
\begin{equation}
 \min_z\sum_{i,j}\Delta_{ij}z_{ij}
 \quad\text{subject to}\quad
 \sum_jz_{ij}=1,\quad\sum_iz_{ij}\le1,\quad z_{ij}\in\{0,1\}.
 \label{eq:assignment-model}
\end{equation}
In this ideal model, disallowed edges are excluded. When the model is feasible,
the Hungarian algorithm finds its optimum \citep{kuhn1955hungarian}.
For the simultaneous placement $p^z$,
define its interaction residual by
\begin{equation}
 R_F(z)=F(p^z)-F(p^0)-\sum_{i,j}\Delta_{ij}z_{ij}.
 \label{eq:assignment-residual}
\end{equation}
The joint placement improves exactly when
$\sum_{ij}\Delta_{ij}z_{ij}+R_F(z)<0$; the residual can have either sign.
For the canonical proxy, linearity of its component weights gives
$R_{\proxy}=R_{\rm wl}+\tfrac12R_{\rm den}+\tfrac12R_{\rm cong}$.
A net touched by at most one moved macro has zero HPWL interaction residual;
shared nets need not. Even disjoint local load changes can interact through
top-$k$ selection in the density and congestion terms.

A small HPWL example shows that the assignment can worsen the joint objective
when all proposed positions are legal. Place two movable macros at
$x_1=2$ and $x_2=12$, at a common vertical coordinate. Each has a zero-offset
pin on a shared weight-two net and a second, coincident pin on a unit-weight
net to a fixed terminal at its initial coordinate. Ignoring the fixed positive
normalization, their HPWL is
\begin{equation}
 F(x_1,x_2)=2|x_1-x_2|+|x_1-2|+|x_2-12|.
 \label{eq:assignment-example}
\end{equation}
For candidate slots $3$ and $11$, the frozen-background cost matrix is
\begin{equation}
 \Delta=\begin{pmatrix}-1&-9\\-9&-1\end{pmatrix}.
\end{equation}
The assignment selects $(x_1,x_2)=(11,3)$ with predicted change $-18$.
Yet $F(2,12)=20$ and $F(11,3)=34$, giving a realized increase of $14$ and
residual $32$. Macro widths below one and a sufficiently large canvas make the
initial, single-move, and final placements non-overlapping. Fixed terminals are
points and impose no macro-overlap constraint. This is an HPWL counterexample,
not a claim about a measured full-proxy run.

The implementation uses the incremental replica for matrix entries. Before
solving, it replaces infinities with finite penalties and sets an all-infinite
row to zero. It then commits moves sequentially, skipping illegal targets.
The realized partial batch can
therefore differ from $p^z$. Its explicit batch-score check measures the
interactions represented by that replica; full rescoring checks the phase
endpoint. Those checks are necessary even with a perfectly solved assignment
and exact individual cost entries.
